\documentclass[12pt]{article}
\usepackage{amsmath,amssymb,amsfonts}
\usepackage[margin=1in]{geometry}

\begin{document}

\begin{center}
{\Large\bfseries Chapter 10, Problem 5}
\end{center}

\noindent\textbf{Problem.} Consider an arbitrary group, $G$, with a subgroup, $H$. Show that if $H$ is of index two, then $H$ is invariant. Remember that the index of a subgroup is the ratio of the order of the group to the order of the subgroup.

\bigskip
\noindent\textbf{Solution.}

Let $|G| = n$ and $|H| = n/2$, so $H$ has index 2 in $G$. This means there are exactly two left cosets of $H$ in $G$: namely $H$ itself and one other coset $gH$ for any $g \notin H$. Similarly, there are exactly two right cosets: $H$ and $Hg$.

We need to show that $gHg^{-1} = H$ for all $g \in G$, i.e., $gH = Hg$ for all $g$.

If $g \in H$, then $gH = H = Hg$, which is trivially true.

If $g \notin H$, consider the left coset $gH$. Since there are only two left cosets, $H$ and $gH$, and these partition $G$, we have:
\[
G = H \cup gH \quad (\text{disjoint union}).
\]

Therefore $gH = G \setminus H$ (the set complement of $H$ in $G$).

Similarly, the two right cosets are $H$ and $Hg$, so:
\[
G = H \cup Hg \quad (\text{disjoint union}),
\]
which gives $Hg = G \setminus H$.

Since both $gH$ and $Hg$ equal $G \setminus H$, we conclude:
\[
gH = Hg \quad \text{for all } g \notin H.
\]

Combined with the trivial case $g \in H$, we have $gH = Hg$ for all $g \in G$. Therefore $H$ is a \textbf{normal (invariant) subgroup} of $G$. $\square$

\end{document}
